\section{Related Work}
\label{sec:related}

\paragraph{Formulaic alphas and automated search.}
\citet{kakushadze2016alphas} documents 101 formulaic alphas built from a small
vocabulary of time-series and cross-sectional operators over daily price and
volume data. The factor language in this paper uses a vocabulary of that kind,
and our reference library contains three of its formulas (translated into
the language) and further entries written in that spirit. Searching
over expression trees is the classic setting for genetic
programming~\citep{koza1992genetic}, which serves as our main non-LLM
baseline. \citet{yu2023alphagen} use reinforcement learning to generate
collections of formulaic alphas that work well together. \citet{wang2023alphagpt}
describe an interactive system in which an LLM helps a human analyst mine
alphas. In our study all proposers share the language, the validator's
limits, the trial budget and the evaluation, so that differences between
methods can be attributed to the proposer rather than to the evaluation. The
baselines, however, sample from a restricted menu of windows, literals and
depths, to which the LLM is not limited (Section~\ref{sec:limitations}).

\paragraph{Machine learning in empirical asset pricing.}
\citet{gu2020empirical} compare machine-learning methods for predicting the
cross-section of returns and find gains from flexible models. Our object is
different. We search over short, human-readable programs whose number of
trials can be counted exactly, so that selection effects can be controlled
formally.

\paragraph{LLM-guided program search.}
\citet{romeraparedes2024funsearch} pair an LLM with a programmatic evaluator
and obtain verifiable mathematical discoveries. In finance the evaluator is
itself noisy and the data cannot be regenerated. This moves the burden of
proof from the proposer to the protocol: multiple-testing control and
out-of-sample discipline.

\paragraph{Multiple testing and data snooping.}
We control the false discovery rate with the step-up procedure of
\citet{benjamini1995fdr}. The step-up procedure of \citet{benjamini2001fdr}
serves as a robustness check under arbitrary dependence. Selection-adjusted
Sharpe-type statistics follow the deflated Sharpe ratio
of~\citet{bailey2014deflated}. \citet{harvey2016cross} document the scale of
the multiple-testing problem in the factor literature, and
\citet{bailey2017pbo} formalize the probability of backtest overfitting. The
bootstrap Reality Check of \citet{white2000reality} and the superior
predictive ability test of \citet{hansen2005spa} test whether the best of
many strategies beats a benchmark. The step-down procedure of
\citet{romano2005stepwise} identifies which strategies do so while
controlling the family-wise error rate. We list these tests as planned
extensions, because they need the full trial ledger that our protocol
records. Inference on mean information
coefficients uses heteroskedasticity- and autocorrelation-consistent standard
errors~\citep{newey1987hac}.

\paragraph{Language models, finance and look-ahead.}
\citet{lopezlira2023chatgpt} study whether ChatGPT's reading of news headlines
predicts stock returns. \citet{glasserman2023lookahead} assess look-ahead bias
in GPT-based return predictions. Domain-specific models such as
BloombergGPT~\citep{wu2023bloomberggpt} are trained on large corpora of
financial text. This shows that published financial evidence can enter
pre-training. LLMs are also known to produce fluent but unsupported
content~\citep{ji2023hallucination}, so we treat the economic rationales
attached to proposals as model claims, not as evidence. Our contamination
diagnostics address look-ahead in three ways:
\begin{itemize}
  \item anonymized prompts;
  \item a pre- versus post-cutoff comparison against a baseline;
  \item value-based (behavioural) novelty against a library of known factors.
\end{itemize}
